
%%%%%%%%%%%%%%%%%%%%%%%%%%%%%%%%%%%%%%
%%   Status e Resultados Esperados  %%
%%%%%%%%%%%%%%%%%%%%%%%%%%%%%%%%%%%%%%

\section[Status]{Status e Resultados Esperados}
\label{sec:resultados}

\subsection{Status do Trabalho}
\label{subsec:status}

\begin{frame}{Status do trabalho}
    \begin{itemize}
        \item[$\checkmark$] Proposta de TCC I redigida: introdução, fundamentação teórica e metodologia.
        \item[$\checkmark$] Ambiente experimental preparado (\textbf{WSL2 Ubuntu}): \textit{Online Boutique} executando sobre o \textit{Service Weaver}.
        \item[$\checkmark$] Revisão da literatura consolidada (11 referências com DOI).
        \item[{\textperiodcentered}] Definição final dos cenários de granularidade (em andamento).
        \item[{\textperiodcentered}] Execução dos experimentos e coleta de métricas (novembro).
    \end{itemize}
\end{frame}

\subsection{Resultados Esperados}
\label{subsec:resultados-esperados}

\begin{frame}{Resultados esperados}
    \begin{itemize}
        \item Consolidar $U(G)$ para cada cenário $\times$ perfil de carga.
        \item \textbf{Hipótese}: o equilíbrio entre escalabilidade e custo \textbf{varia com o perfil de carga}.
        \item Evidências para orientar decisões de granularidade.
    \end{itemize}
    \begin{center}
    \begin{tikzpicture}[font=\scriptsize, scale=0.8]
        % eixos
        \draw[-{Latex[length=2mm]}] (0,0) -- (5.0,0) node[below left=1pt]{cenários};
        \draw[-{Latex[length=2mm]}] (0,0) -- (0,2.9) node[left, rotate=90, anchor=south]{utilidade $U(G)$};
        % barras: carga moderada (azul)
        \fill[Lightblueframe, draw=Blueframe!70!black] (0.62,0)   rectangle (0.98,1.2);
        \fill[Lightblueframe, draw=Blueframe!70!black] (1.82,0)   rectangle (2.18,1.9);
        \fill[Lightblueframe, draw=Blueframe!70!black] (3.02,0)   rectangle (3.38,1.5);
        % barras: pico (laranja)
        \fill[Lightorangeframe, draw=Orangeframe!70!black] (1.02,0) rectangle (1.38,0.7);
        \fill[Lightorangeframe, draw=Orangeframe!70!black] (2.22,0) rectangle (2.58,1.4);
        \fill[Lightorangeframe, draw=Orangeframe!70!black] (3.42,0) rectangle (3.78,2.1);
        % rótulos dos cenários
        \node at (1.0,-0.28) {A};
        \node at (2.2,-0.28) {B};
        \node at (3.4,-0.28) {C};
        % legenda
        \fill[Lightblueframe, draw=Blueframe!70!black]   (4.35,2.30) rectangle +(0.26,0.26);
        \node[anchor=west, font=\tiny] at (4.68,2.43) {carga moderada};
        \fill[Lightorangeframe, draw=Orangeframe!70!black] (4.35,1.90) rectangle +(0.26,0.26);
        \node[anchor=west, font=\tiny] at (4.68,2.03) {pico de acesso};
    \end{tikzpicture}
    \end{center}
    \vspace{-2mm}
    \centering\scriptsize Ilustração do formato esperado dos resultados (valores fictícios). Cenários: A = monólito, B = intermediária, C = fina.
\end{frame}
